\clearpage
\fancyhead[L]{\small IRIS.}
\SourceLeaf{478}
\begin{center}
\rule{0.55\linewidth}{0.4pt}\\[2pt]
Nachrichten der Königl. Gesellsch. d. Wiss. zu Göttingen. Nr. 11. 13. September 1847.\\[2pt]
\rule{0.55\linewidth}{0.4pt}
\end{center}

Von der Entdeckung eines neuen Planeten durch Herrn \textsc{Hind} am 13. August in London erfuhr ich zuerst aus einer Pariser Zeitung vom 16. August. Aufgrund dieser Meldung beobachtete ich ihn am 21. August im Meridian. Da die Meldung nur zwei Positionen im Abstand von einer Stunde enthielt, blieb zunächst eine gewisse Unsicherheit über seine Identität. Bald darauf wurde sie nach Eingang weiterer englischer Beobachtungen vom 14. und 15. August ausgeräumt. Seitdem habe ich den Planeten, der bekanntlich den Namen Iris erhalten hat, an jedem Abend, an dem es das Wetter erlaubte, am Meridiankreis beobachtet. Bisher habe ich folgende Positionen erhalten.

\begin{center}
\setlength{\tabcolsep}{4pt}
\begin{tabular}{lr|r|r|r}
\multicolumn{2}{c|}{1847} & \shortstack{Mittlere Zeit\\in Göttingen} & Rektaszension & \shortstack{Südliche\\Deklination} \\\hline
August & 21 & $9^{\rm u}\ 52^{\rm m}\ 54^{\rm s}1$ & $297^\circ\ 47'\ 20''8$ & $13^\circ\ 42'\ 19''9$ \\
 & 28 & $9\ 21\ 16.6$ & $296\ 45\ 38.2$ & $13\ 54\ 47.9$ \\
 & 29 & $9\ 16\ 53.3$ & $296\ 38\ 46.6$ & $13\ 56\ 34.2$ \\
Sept. & 4 & $8\ 51\ 12.0$ & $296\ 7\ 14.1$ & $14\ 6\ 3.8$ \\
 & 6 & $8\ 42\ 54.5$ & $296\ 0\ 46.5$ & $14\ 8\ 56.1$ \\
\end{tabular}
\end{center}

Bei diesen Beobachtungen (unter meist wenig günstigen atmosphärischen Bedingungen) erschien der Planet nur so hell wie ein Stern 9. Größe.

Sobald die Beobachtung vom 28. August reduziert war, begann Herr Professor \textsc{Goldschmidt} mit der Berechnung der elliptischen Elemente. Er legte meine beiden Beobachtungen vom 21. und 28. August sowie eine Beobachtung zugrunde, die Herr \textsc{Challis} am 14. August in Cambridge gemacht hatte. Das sind seine Ergebnisse.

\begin{center}
\small
\setlength{\tabcolsep}{2.5pt}
\begin{tabular}{llr}
\multicolumn{2}{l}{\shortstack[l]{Mittlere Anomalie 1847 August 14. $0^{\rm u}$,\\Mittl. Göttinger Zeit}} & $288^\circ\ 53'\ 21''9$ \\
Länge des Perihels & \mbox{.\ .\ .\ .\ .\ .\ .\ .\ .\ .\ .\ .\ .\ .} & $44\ 6\ 49.3$ \\
Länge des aufsteigenden Knoten & \mbox{.\ .\ .\ .\ .\ .\ .\ .\ .} & $260\ 27\ 58.8$ \\
\multicolumn{3}{l}{beide vom scheinbaren Äquinoktium des 14. August an gerechnet} \\
Mittlere tägliche siderische Bewegung & \mbox{.\ .\ .\ .\ .\ .\ .} & $982''5907$ \\
Logarithmus der großen Halbachse & \mbox{.\ .\ .\ .\ .\ .\ .\ .} & $0.3717560$ \\
Exzentrizität & \mbox{.\ .\ .\ .\ .\ .\ .\ .\ .\ .\ .\ .\ .\ .\ .\ .} & $0.2135543$ \\
Neigung der Bahn & \mbox{.\ .\ .\ .\ .\ .\ .\ .\ .\ .\ .\ .\ .\ .} & $5^\circ\ 29'\ 55''8$ \\
\end{tabular}
\end{center}

Herr \textsc{Gould} aus Boston, der sich bei uns mit großem Eifer astronomischen Studien widmet, hat aus denselben Beobachtungen ebenfalls eine elliptische Bahn berechnet. Sie unterscheidet sich nur wenig von der oben angegebenen und stimmt auch mit den späteren, bis jetzt gemachten Beobachtungen noch sehr gut überein.

Zwei Schlüsse lassen sich daraus bereits ziehen. Erstens steht nun fest, dass Iris ein Planet ist und dass seine Bahn ebenso wie die Bahnen von Ceres, Pallas, Juno, Vesta, Asträa und Hebe zwischen den Bahnen von Mars und Jupiter liegt. Zweitens wird der Planet bei seiner Erscheinung im Jahr 1848 voraussichtlich deutlich heller sein als in diesem Jahr.

\noindent Göttingen, den 7. September 1847.\hfill\textsc{Gauss}.
